\documentclass[11pt,a4paper]{article}
\usepackage{amsmath,amssymb}
\usepackage[margin=2.5cm]{geometry}
\usepackage{array}

\title{Mamba-3 SSD: Explicit Attention Documentation}
\author{}
\date{}

\begin{document}
\maketitle

\section{Core Definitions \& Activations}

Before projections, the following mathematical functions and operators are defined:

\begin{itemize}
  \item \textbf{Softplus:} $\text{Softplus}(x) = \ln(1 + \exp(x))$. This function is applied to the step-size projection to ensure $\Delta_i > 0$, which is a mechanical necessity for the stability of the state-space decay.
  \item \textbf{SiLU (Swish):} $\text{SiLU}(x) = x \cdot \sigma(x)$. This is used for the value projection $v_i$ to provide non-linear gating.
  \item \textbf{Hadamard Product ($\odot$):} Element-wise multiplication.
  \item \textbf{Matrix Dualism:} In the SSD framework, \textbf{Self-Attention} is the parallel matrix representation of the recurrent state evolution.
\end{itemize}

\section{Parameter Dimensions Reference}

To maintain consistency in your implementation, all variables and learned weight matrices follow these dimensions:

\begin{center}
\begin{tabular}{|l|l|l|}
\hline
\textbf{Variable} & \textbf{Dimension} & \textbf{Definition} \\
\hline
$x_i$ & $D \times 1$ & Input feature vector for patch $i$ ($L = T \times H \times W$). \\
\hline
$W_B, W_C$ & $N \times D$ & \textbf{Key} and \textbf{Query} projection matrices ($N$ = state dimension). \\
\hline
$W_{\Delta}$ & $1 \times D$ & \textbf{Step-size (delta)} projection row vector. \\
\hline
$W_V$ & $D \times D$ & \textbf{Value} projection matrix. \\
\hline
$W_G$ & $N \times D$ & \textbf{Cross-attention gating} projection matrix. \\
\hline
$A$ & $1 \times 1$ & Learned state-space evolution scalar. \\
\hline
$h_i$ & $N \times D$ & Hidden state matrix at step $i$. \\
\hline
$y_i$ & $D \times 1$ & Final output feature vector. \\
\hline
\end{tabular}
\end{center}

\section{Global Self-Attention (Intra-Sequence Context)}

This mechanism allows every patch in the image sequence to interact with every other patch via a symmetric distance decay.

\subsection{Input Projections}

For each patch token $x_i$:

\begin{enumerate}
  \item \textbf{Key ($B_i$):} $B_i = W_B x_i \in \mathbb{R}^{N \times 1}$
  \item \textbf{Query ($C_i$):} $C_i = W_C x_i \in \mathbb{R}^{N \times 1}$
  \item \textbf{Step Size ($\Delta_i$):} $\Delta_i = \text{Softplus}(W_{\Delta} x_i) \in \mathbb{R}^{1 \times 1}$
  \item \textbf{Value ($v_i$):} $v_i = \text{SiLU}(W_V x_i) \in \mathbb{R}^{D \times 1}$
\end{enumerate}

\subsection{Interaction Equation}

The output $y_{self, i}$ is the result of the Query $C_i$ interacting with the global weighted sum of all Key-Value pairs:

\begin{equation}
y_{self, i} = \sum_{j=1}^L \mathcal{L}_{ij} \cdot (C_i^\top B_j) v_j
\end{equation}

\textbf{Symmetric Distance Decay ($\mathcal{L}_{ij}$):}

\begin{equation}
\mathcal{L}_{ij} = \exp\left( - \left| \sum_{k=\min(i,j)}^{\max(i,j)} \Delta_k A \right| \right)
\end{equation}

\section{Cross-Attention (Multi-View/Reference Context)}

Cross-attention facilitates the interaction between the target image sequence and a pre-computed reference view state.

\subsection{Context State Acquisition}

The reference view sequence ($L_{ref}$) is compressed into a static state matrix $h_{ctx} \in \mathbb{R}^{N \times D}$:

\begin{equation}
h_{ctx} = \sum_{j=1}^{L_{ref}} \left( \prod_{k=j+1}^{L_{ref}} \alpha_{k, ref} \right) B_{j, ref} v_{j, ref}^\top
\end{equation}

\textit{(Where $\alpha_k = \exp(-\Delta_k A)$)}

\subsection{Interaction Equation}

For each target token $x_{target, i}$, the output $y_{cross, i}$ queries the combined local and reference context:

\begin{equation}
y_{cross, i} = C_{target, i}^\top \left[ h_{local, i} + (\text{Gate}_i \odot h_{ctx}) \right]
\end{equation}

\begin{itemize}
  \item $h_{local, i}$: The current recurrent state of the target sequence.
  \item $\text{Gate}_i$: $\sigma(W_G x_{target, i}) \in \mathbb{R}^{N \times 1}$, a data-dependent filter for the reference state.
\end{itemize}

\section*{Implementation Note}

For your goal of surpassing SOTA in 3D reconstruction by March 2026, the use of \textbf{Trapezoidal Discretization} in the $\mathcal{L}_{ij}$ matrix is essential. It ensures that the model captures the fine-grained geometric gradients (surface normals and curvatures) better than the standard Euler-based Mamba-2.

\end{document}
